\subsection{The Traditional Compilation Pipeline (The C Blueprint)}

\subsubsection{Core Definition, Intent, and Objectives of Compilation vs. Interpretation}

Compilation transforms a human program into another representation \emph{before} execution; interpretation keeps a runtime mediator active while the program runs. C's classic path ends in a native executable the operating system loads directly. Python's classic path loads \texttt{python} first; the \texttt{.py} file becomes input to CPython, which compiles to bytecode and interprets it inside an already-running process.

The architectural question is therefore not ``compiled or interpreted?'' but \emph{where} translation, validation, optimization, and binding occur. C pushes most binding and layout decisions to build time. Python defers type and object layout to runtime, trading early machine specificity for flexibility inside a managed engine.

\subsubsection{The Factory vs. The Traveling Theater}

Think of the traditional C toolchain as an \textbf{industrial factory} and a later Python deployment as a \textbf{traveling theater} that erects its stage inside an existing building.

\paragraph{Stage 1: Raw materials intake (source and preprocessing).}
The factory receives blueprints (translation units). Preprocessing is the clerical desk: macros expand, headers declare contracts, and conditional compilation selects which blueprint sheets enter production. Nothing executes yet; the factory only normalizes what the assembly line will see.

\paragraph{Stage 2: Token stamping (lexical analysis).}
A scanner does not ``understand'' programs. It stamps uniform tokens---identifiers, literals, operators, punctuation---onto the conveyor. In our later Python story, this stage still exists, but Python's lexer also emits \texttt{INDENT} and \texttt{DEDENT} tokens from physical layout, a design choice C's factory deliberately ignores.

\paragraph{Stage 3: Structural framing (parsing and the AST).}
The parser validates grammar and erects a \textbf{wooden scale model}: the abstract syntax tree. The model is target-independent structure: expressions, statements, declarations. It is precise enough to inspect and rewrite, but not yet the final product.

\paragraph{Stage 4: Safety inspection (semantic analysis).}
Engineers attach meaning: symbol tables, types, scopes, linkage, and illegal combinations are rejected \emph{before} concrete is poured. This is where C earns its reputation for early failure: many mistakes never reach the CPU.

\paragraph{Stage 5: Industrial optimization and lowering (middle/back end).}
The factory reshapes the model for one soil type: a specific ISA and ABI. Optimizations rewrite intermediate representations; instruction selection and register allocation choose how the machine will actually move bits. The output is not Python bytecode; it is assembly or object code destined for a particular CPU family.

\paragraph{Stage 6: Permanent pour (assembly, linking, native executable).}
Object files are cast modules. The linker pours them into one loadable executable: machine code in the text segment, initialized data, BSS, relocation metadata, and an entry point. When the factory leaves, only the concrete remains. The original blueprints are not consulted at runtime.

\paragraph{Why this metaphor matters for Python.}
CPython repeats the \emph{same structural stations} with different materials:
\begin{itemize}
    \item tokens and parsing still build an AST;
    \item semantic passes still attach meaning, but many rules are deferred to runtime objects;
    \item the ``pour'' is not ISA-specific machine code for the user program, but portable bytecode plus runtime support inside the interpreter binary.
\end{itemize}

The traveling theater does not demolish the host building (the OS process running \texttt{python}). It brings its own stage machinery: bytecode interpreter loop, object heap, frames, and namespaces. Performances (user code) can change nightly; the concrete shell (CPython itself) was poured earlier by a C factory.

\paragraph{Compilation objectives restated compactly.}
\begin{itemize}
    \item \textbf{translation}: high-level structure $\rightarrow$ lower-level executable representation;
    \item \textbf{validation}: reject ill-formed or ill-typed programs early;
    \item \textbf{optimization}: rewrite for speed, size, and instruction-level efficiency;
    \item \textbf{target binding}: specialize to one processor/OS executable format;
    \item \textbf{composition}: link multiple translation units and libraries;
    \item \textbf{early layout}: fix many offsets and call conventions before execution begins.
\end{itemize}

For a C programmer moving to Python, the decisive shift is ownership of execution: in C, the user's compiled image \emph{is} the process image; in Python, the interpreter process owns execution and the user's module is data plus bytecode consumed inside that process.
